\documentclass[10pt,a4paper]{amsart}
\usepackage[margin=19mm]{geometry}
\usepackage{amsmath,amssymb,mathtools}
\usepackage[scr=rsfs,cal=euler]{mathalfa}
\usepackage{xfrac}
\usepackage[unicode,hidelinks,bookmarks=false]{hyperref}
\newcommand{\ie}{i.e.\ }
\newcommand{\cf}{cf.\ }
\newcommand{\resp}{respectively\ }
\newcommand{\Subsection}{\subsection}
\providecommand{\nobreakdash}{\nobreak\hbox{-}}
\setlength{\emergencystretch}{2em}
\multlinegap0pt
\usepackage{bm}
\usepackage{bbm}
\usepackage{dsfont}
\usepackage{xspace}
\usepackage{cancel}
\usepackage{mathtools}
\usepackage{cleveref}
\usepackage{amsthm}
\makeatletter
\@ifundefined{Proposition}{\newtheorem{Proposition}{Proposition}}{}
\@ifundefined{Lemma}{\newtheorem{Lemma}[Proposition]{Lemma}}{}
\makeatother
\DeclareMathOperator{\Aut}{Aut}
\DeclareMathOperator{\Sym}{Sym}
\title[Branch iterated Galois groups with positive fixed-point prop]{Branch iterated Galois groups with positive fixed-point proportion and positive Hausdorff dimension}
\author{Santiago Radi}
\date{Source: arXiv:2602.13428v3 (CC BY 4.0)}
\begin{document}
\maketitle

\noindent\emph{Source and licence.} Branch iterated Galois groups with positive fixed-point proportion and positive Hausdorff dimension. Version arXiv:2602.13428v3 (CC BY 4.0), distributed under the Creative Commons Attribution 4.0 International licence, as recorded in the arXiv licence metadata for this exact version and shown on the abstract page https://arxiv.org/abs/2602.13428v3. Full rights notice: licenses/arxiv-2602.13428-CC-BY-4.0.txt.\par\smallskip

\noindent\textit{Editorial scope.} Counted unit: the Lemma whose source label is \texttt{lemma: GQP group and index} in the article (source0.tex lines 686--694, source label \texttt{lemma: GQP group and index}), delivered together with the article's own complete original proof of it (source0.tex lines 696--736, one \texttt{proof} environment of 41 lines). No mathematics is written, repaired or re-derived: statement and proof are the article's own text line for line. Required context retained uncounted: equation: properties sections (lines 237--240). Every definition, display and label of those lines is retained unchanged. Omitted from this package and not redistributed: 1--236, 241--685, 695--695, 737--1187. Nothing inside a retained excerpt is omitted or altered: every retained body is delivered exactly as the article's lines are written, so no omission receipt of any kind accompanies this package. Citations: the retained excerpts carry no citation command at all, so no bibliography entry of the article is transcribed and no reference is redistributed. Reference adaptation: none was needed and none was made; every reference inside the delivered unit resolves inside it, so no unresolved reference or citation is produced. Printed slips kept literal: no printed word, symbol, hypothesis or number of the counted statement or of its proof was altered. Layout and class: the article's own class is \texttt{amsart} with packages including inputenc, fontenc, textcomp, amsfonts, amsthm, amssymb, pst-node, xy, mathtools, chngcntr, thmtools, mathrsfs, tcolorbox, listings; the wrapper is the lane's pinned \texttt{amsart} editorial wrapper at 10pt on A4, which restates the theorem-like environments the retained text needs and adds the article's own macro definitions that the retained text needs (\texttt{\textbackslash Aut}, \texttt{\textbackslash Sym}), with extra wrapper packages (bm, bbm, dsfont, xspace, cancel, mathtools, cleveref). The delivered numbering is the unit's own. Assets: no retained span contains a figure environment, a graphics-inclusion command, a PDF-page inclusion, a table or a TikZ picture; no image file is copied into this package, the package declares no asset dependency and no image file is redistributed with it. Licence: version arXiv:2602.13428v3 of this article is distributed under the Creative Commons Attribution 4.0 International licence, as recorded in the arXiv licence metadata for this exact version and shown on the abstract page https://arxiv.org/abs/2602.13428v3. A full rights notice is delivered at licenses/arxiv-2602.13428-CC-BY-4.0.txt. Characters: every retained excerpt is pure ASCII, so the delivered engine, XeLaTeX with its default Latin Modern text font, reports no missing character.
\begin{align}
(gh)|_v^n = g|_{h(v)}^n h|_v^n \text{ and } (g^{-1})|_v^n = (g|_{g^{-1}(v)}^n)^{-1}.
\label{equation: properties sections}
\end{align}

\begin{Lemma}
Let $T$ be a $d$-regular tree with $d \geq 2$ and $1 \neq \mathcal{Q} \lhd \mathcal{P} \leq \Sym(d)$. Then, the group $G_\mathcal{Q}^\mathcal{P}$ satisfies the following properties:
\begin{enumerate}
\item $G_\mathcal{Q}^\mathcal{P}$ contains $W_\mathcal{Q}$,
\item $W_\mathcal{Q} \lhd G_\mathcal{Q}^\mathcal{P}$,
\item $[G_\mathcal{Q}^\mathcal{P} : W_\mathcal{Q}] = [\mathcal{P}: \mathcal{Q}]$.
\end{enumerate}
\label{lemma: GQP group and index}
\end{Lemma}

\begin{proof}
We start proving that $G_\mathcal{Q}^\mathcal{P}$ is a group. Clearly, the identity element is in $G_\mathcal{Q}^\mathcal{P}$. Next, suppose $g,h \in G_\mathcal{Q}^\mathcal{P}$. We must verify that $gh^{-1} \in G_\mathcal{Q}^\mathcal{P}$. For $v,w \in T$, by the properties of sections (see \cref{equation: properties sections}), we have:
\begin{align*}
(gh^{-1})|_v^1 \, \, ((gh^{-1})|_w^1))^{-1} = (g|_{h^{-1}(v)}^1) \, (h|_{h^{-1}(v)}^1)^{-1} \, (h|_{h^{-1}(w)}^1) \,  (g|_{h^{-1}(w)}^1)^{-1}.
\end{align*}

Rewriting, this becomes:

\begin{align*}
\left[ (g|_{h^{-1}(v)}^1) \, (h|_{h^{-1}(v)}^1)^{-1} \, (h|_{h^{-1}(w)}^1) \, (g|_{h^{-1}(v)}^1)^{-1} \right]  \left[ ((g|_{h^{-1}(v)}^1)  (g|_{h^{-1}(w)}^1)^{-1}) \right].
\end{align*}

The second term between brackets is in $\mathcal{Q}$ since $g \in G_\mathcal{Q}^\mathcal{P}$ and the first term in brackets is in $\mathcal{Q}$ because $h \in G_\mathcal{Q}^\mathcal{P}$ and $\mathcal{Q} \lhd \mathcal{P}$. This ensures that $gh^{-1} \in G_\mathcal{Q}^\mathcal{P}$. Hence, $G_\mathcal{Q}^\mathcal{P}$ is a group.

We now prove that $W_\mathcal{Q} \lhd G_\mathcal{Q}^\mathcal{P}$. The proof is similar to the previous calculation. Specifically, we need to show that if $g \in G_\mathcal{Q}^\mathcal{P}$ and $h \in W_\mathcal{Q}$ then $(ghg^{-1})|_v^1 \in \mathcal{Q}$ for all $v \in T$. Indeed,
\begin{align*}
(ghg^{-1})|_v^1 = g|_{hg^{-1}(v)}^1 \, h|_{g^{-1}(v)}^1 \, (g|_{g^{-1}(v)}^1)^{-1}.
\end{align*}

Rewriting, this becomes:

\begin{align*}
 (ghg^{-1})|_v^1 = \left[ g|_{hg^{-1}(v)}^1 \, h|_{g^{-1}(v)}^1 \, (g|_{hg^{-1}(v)}^1)^{-1} \right] \, \left[ (g|_{hg^{-1}(v)}^1) \, (g|_{g^{-1}(v)}^1)^{-1} \right].  
\end{align*}

The first term in brackets belongs to $\mathcal{Q}$ because $\mathcal{Q} \lhd \mathcal{P}$, and the second term in bracket belongs to $\mathcal{Q}$ since $g \in G_\mathcal{Q}^\mathcal{P}$. Thus, $(ghg^{-1})|_v^1 \in \mathcal{Q}$ and consequently $W_\mathcal{Q} \lhd G_\mathcal{Q}^\mathcal{P}$.

Finally, consider the map 
\begin{align}
\begin{split}
G_\mathcal{Q}^\mathcal{P}/W_\mathcal{Q} & \rightarrow \mathcal{P}/\mathcal{Q} \\ 
[g] & \mapsto [g|_\emptyset^1].
\label{equation: cosets of GQP}
\end{split}
\end{align}
The map is well-defined because if $h \in W_\mathcal{Q}$ and $g \in G_\mathcal{Q}^\mathcal{P}$ then $$[(gh)|_\emptyset^1] = [g|_\emptyset^1] [h|_\emptyset^1] = [g|_\emptyset^1],$$ and it is in fact a morphism since the root of the tree is fixed by any element. 

To show surjectivity, let $\sigma \in \mathcal{P}$.  Define $g \in \Aut(T)$ such that $g|_v^1 = \sigma$ for all $v \in T$. Then, $g \in G_\mathcal{Q}^\mathcal{P}$ and $[g|_\emptyset^1] = [\sigma]$. 

To prove injectivity, suppose $[g|_\emptyset^1] = [h|_\emptyset^1]$. Then: $$(g|_\emptyset^1)(h|_\emptyset^1)^{-1} = (gh^{-1})|_\emptyset^1 \in \mathcal{Q}.$$ Since $gh^{-1} \in G_\mathcal{Q}^\mathcal{P}$, then $(gh^{-1})|_v^1 \in \mathcal{Q}$ for all $v \in T$ and consequently $gh^{-1} \in W_\mathcal{Q}$. 
\end{proof}

\end{document}
